% TRANSLATION-PRODUCER DRAFT — UNCHECKED
% Work: Emmy Noether, Paper 41 — Korean U04
% Source authority custody: ${PUBLIC_INTERLANGUAGE_ROOT}\03_projects\language_management\cjk\03_working_translations\noether_paper41_zh_translation_001_20260722\source\Noether_Paper41_CurrentGermanAuthority_interval.tex
% Exact source snapshot locator: lines 48--56
% Source unit LF-normalized UTF-8 SHA-256: F0D1FF2818BD95257249E93DB56E891343DF2EB025DD3F46987461ACCA05A0FC
% Source snapshot SHA-256: C265058425E5E2D1A2289CC03A9DDEDDDF4803A3215DC3F173B93E7AB69D60ED
% Historical whole-source SHA-256: 443EF950D7D45DC6D9E44A9B87501620C10DA873E50E5F2B253ECCAE6A946D27
% Role boundary: translation only. No source/scan check, Korean review or self-review, compilation, rendering, assembly, packaging, certification, approval, or German-source alteration was performed.
% Checker handoff state: UNCHECKED; all Korean wording and preserved TeX require independent checking.

이제 순환대수인 특수 경우를 간단히 제시한다. \(Z/k\)가 순환이고 \(S\)가 그 군의 생성 치환이라 하자. 그러면 \(S\)의 거듭제곱들에 \(u\)의 거듭제곱들을 대응시킬 수 있으며, 다음을 얻는다.

\srcdisplay{(1')}{A=Z+uZ+\cdots+u^{n-1}Z.}
\srcdisplay{(2')}{zu=uz^S.}
\srcdisplay{(3')}{u^n=\alpha.}
\srcdisplay{(4')}{\alpha\text{는 기초체 }k^*\text{에 속한다}.}
\srcdisplay{(5')}{\bar\alpha=\alpha\cdot N(c),\quad\text{단 }v=uc\text{로 둔다}.}

따라서 여기에서 각 인자계는 기초체에 속하는 단 하나의 원소 \(\alpha\)로 이루어지며---\(A=(\alpha,Z,S)\)로 표기한다---인자계의 단위원류는 \(Z^*\)에서 오는 노름들로 주어지고, 대수류군은 몫군 \(k^*/N(Z^*)\)와 동형이 된다.
